\section{Introduction}

Full-model sliding window attention (SWA) conversion of pre-trained transformers fails catastrophically without fine-tuning \cite{Yu2025SWAA}, yet per-layer analysis of Llama-2-7B reveals that attention weight is already concentrated in fewer than 10\% of tokens (Gini coefficient $= 0.6829$, top-10\% token share $= 71.72\%$)---a near-universal structural property confirmed across 200 diverse evaluation examples with zero exceptions. This tension raises a precise question: if many layers are already operating with effectively local attention patterns, why must fine-tuning be the only path to safe SWA conversion?

The practical motivation is substantial. Transformer attention scales as $O(n^2)$ per layer, making long-context inference expensive. SWA reduces this to $O(n \cdot w)$ per layer by restricting each query to attend only within a local window of width $w$. In production, Mistral 7B uses window-based attention trained from scratch \cite{Jiang2023Mistral}, and Longformer's sliding window achieves 5,840+ citations as a foundational efficiency primitive \cite{Beltagy2020Longformer}. However, \emph{post-hoc} conversion of an already-trained model to SWA requires either fine-tuning recovery \cite{Yu2025SWAA,Liu2026Architecture} or accepting severe quality degradation. A zero-shot, training-free selective conversion criterion would enable deployment-time efficiency gains without the compute overhead of re-training.

The deeper problem is that existing approaches treat the model holistically: either all layers keep full attention or the model is fine-tuned after structural changes. This ignores a critical source of heterogeneity---individual layers differ substantially in how diffuse or concentrated their attention patterns are. Head-removal studies \cite{Michel2019Sixteen} and encoder analysis \cite{Raganato2020Fixed} establish that many attention heads are removable or fixed to positional patterns, but these insights have not been translated into a principled, layer-level, zero-shot criterion for SWA conversion in causal decoder LLMs.

The gap is clear: no method provides a calibration-only signal for identifying which specific layers in a pre-trained decoder LLM can be safely converted to SWA without fine-tuning and without quality degradation exceeding 2 perplexity points.

\textbf{Our key insight:} Head-mean attention entropy, computed over a small calibration set (100 sequences, single GPU, $\approx 20$--30 GPU-minutes for a single calibration pass), provides a stable, deterministic ranking of Llama-2-7B layers by their attention concentration structure. Layers with \emph{high} entropy have diffuse attention---they already spread weight broadly across many tokens without sharp global retrieval. These are the layers where constraining attention to a local window removes behavior they were not strongly exploiting. Critically, the aggregation method matters: head-mean pooling yields Gini $= 0.681$, while head-max pooling yields Gini $= 0.466$---a 32\% relative reduction---because head-max is dominated by individual outlier heads rather than capturing layer-level structure.

\textbf{Paper scope:} This paper is a complete, confirmatory prerequisite validation study. We fully validate the entropy criterion itself (h-e1): concentration characterization, pooling method comparison, and ranking stability. The accuracy-preservation validation (h-e2, h-m1, h-m2) is the planned follow-up, described here as the natural scientific continuation. Readers should understand this paper's confirmed contribution as establishing that a \emph{stable, calibration-only entropy criterion exists}---the question of whether that criterion successfully identifies SWA-safe layers is the next experimental step.

This paper makes the following contributions:

\textbf{(1) Attention Concentration Characterization (Empirical, Confirmed):} We establish that Llama-2-7B exhibits strong, near-universal within-layer attention concentration: Gini mean $= 0.6829$ (std $= 0.0117$), top-10\% token share mean $= 0.7172$ (std $= 0.0196$), with 100\% of 200 evaluation examples satisfying both criteria simultaneously. This is the first systematic application of per-layer head-mean entropy pooling to characterize attention concentration in Llama-2-7B, specifically targeting layer selection for SWA conversion.

\textbf{(2) Pooling Method as a First-Class Design Choice (Methodological, Confirmed):} We demonstrate that head-mean pooling is substantially more stable than head-max for layer-level entropy characterization, producing 32\% higher measured Gini concentration ($0.681$ vs $0.466$). This elevates pooling aggregation from an implementation detail to a methodological decision with measurable impact on signal quality.

\textbf{(3) Entropy Criterion Stability (Empirical, Confirmed):} Layer entropy rankings are stable across calibration subsets, with Spearman $\rho \geq 0.8$ confirmed across all subset pairs, ensuring that the selection criterion is deterministic and not sensitive to the specific 100-sequence subset used.

\noindent\textit{The following contribution describes the designed framework, pending h-e2 confirmation:}

\textbf{(4) Entropy-Guided Selective SWA Framework (Design Only---Experimental Validation Pending):} We design a zero-shot framework for selective SWA conversion of Llama-2-7B, converting the $k = 4$ highest-entropy layers to SWA($w = 512$). The design is complete and implemented; whether this conversion maintains WikiText-103 perplexity within 2 points of the full-attention baseline (our primary prediction P1) is the central open question, requiring h-e2 experimental execution. This contribution documents the framework design, not a validated outcome.

We organize the paper as follows. Section~\ref{sec:related} situates the work within attention efficiency, entropy-based analysis, and model modification literature. Section~\ref{sec:methodology} describes the entropy scoring methodology and selective SWA conversion framework. Section~\ref{sec:experiments} presents experimental design. Section~\ref{sec:results} reports results from the confirmed entropy characterization experiments (h-e1), and describes the planned experimental structure for follow-up work. Section~\ref{sec:discussion} discusses implications and limitations. Section~\ref{sec:conclusion} concludes with the validated prerequisite and open questions for full hypothesis confirmation.
